% BHU PYQ -- Manually transcribed & error-corrected
% Paper: BCF-213 : Corporate Taxation
% Exam: B.Com. (Hons.) F.M.M. Semester III Examination, 2016-17

\documentclass[11pt,a4paper]{article}
\usepackage[a4paper,top=2.5cm,bottom=2.5cm,left=2.8cm,right=2.8cm,headheight=14pt]{geometry}
\usepackage{amsmath,amssymb}
\usepackage[T1]{fontenc}
\usepackage[utf8]{inputenc}
\usepackage{lmodern,microtype,fancyhdr,enumitem,hyperref,lastpage,parskip,booktabs}

\hypersetup{
    colorlinks=true,
    urlcolor=black,
    linkcolor=black,
    pdftitle={BCF-213: Corporate Taxation -- BHU B.Com. (Hons.) FMM Sem III 2016-17}
}

\pagestyle{fancy}
\fancyhf{}
\renewcommand{\headrulewidth}{0.4pt}
\renewcommand{\footrulewidth}{0.4pt}
\fancyhead[L]{\small   \textit{Banaras Hindu University}}
\fancyhead[R]{\small   \textit{BCF-213: Corporate Taxation}}
\fancyfoot[C]{\small Page  \thepage\ of \pageref{LastPage}}


\newlist{parts}{enumerate}{2}
\setlist[parts,1]{label=(\alph*),leftmargin=*,itemsep=5pt,topsep=3pt}
\setlist[parts,2]{label=(\roman*),leftmargin=*,itemsep=3pt,topsep=2pt}

\newcommand{\pts}[1]{\hfill   \textit{[#1]}}


\begin{document}


\begin{center}
  {\large\bfseries BANARAS HINDU UNIVERSITY}\\[2pt]
  {
\normalsize B.Com. (Hons.) F.M.M. --- Semester III Examination, 2016-17}\\[6pt]
  \rule{0.8\linewidth}{0.5pt}\\[6pt]
  {\large\bfseries Paper No. BCF-213: Corporate Taxation}\\[4pt]
  {
\normalsize Time: 3 Hours\hfill Maximum Marks: 70}
\end{center}

\rule{\linewidth}{0.4pt}

\smallskip



\noindent   \textit{Instructions: Answer all questions. All questions carry equal marks (14 marks each).}

\bigskip




\noindent   \textbf{Question 1.}
``Income tax is a tax on income and not receipts.'' Discuss this statement and give the essential characteristics of income.\pts{14}

\centerline{   \textbf{OR}}

Write notes on the following:\pts{14}

\begin{parts}
  \item Assessee
  \item Gross Total Income vs Total Income
  \item Previous year and its exceptions
\end{parts}

\medskip\rule{\linewidth}{0.2pt}




\noindent   \textbf{Question 2.}
Give the name of any fourteen exempted incomes.\pts{14}

\centerline{   \textbf{OR}}

From the following Profit and Loss Account of Mr. Agrawal for the year ended 31st March, 2016, compute his income from business:\pts{14}

\begin{center}
  
\begin{tabular}{p{4.1cm}r|p{4.1cm}r}
     oprule
   \textbf{Debit} &   \textbf{Amount (Rs.)} &   \textbf{Credit} &   \textbf{Amount (Rs.)} \\
    \midrule
    Trade expenses & 150 & Gross Profit & 2,35,500 \\
    Establishment charges & 2,200 & Dividend from domestic company & 2,600 \\
    Rent, rates and taxes & 1,400 & Rent from property & 900 \\
    Income Tax & 700 & Bad debts recovered (earlier disallowed) & 2,000 \\
    Postage and telegram & 100 & & \\
    Gifts and presents & 500 & & \\
    Publicity & 1,575 & & \\
    Fire Insurance Premium & 250 & & \\
    Charity & 375 & & \\
    Donation & 400 & & \\
    Repairs and renewals & 250 & & \\
    Audit fees & 250 & & \\
    Net Profit (transferred to Capital A/c) & 2,32,850 & & \\
    \midrule
   \textbf{Total} &   \textbf{2,41,000} &   \textbf{Total} &   \textbf{2,41,000} \\
    ottomrule
  \end{tabular}
\end{center}

\medskip\rule{\linewidth}{0.2pt}




\noindent   \textbf{Question 3.}
A, B and C are partners sharing profit and loss equally in a firm. The Profit and Loss Account for the year ended 31st March, 2016 showed a net profit of Rs.\ 1,80,000 after debiting Rs.\ 7,000 for interest paid to A @ 20\%, Rs.\ 60,000 for salary paid to B, and Rs.\ 44,000 paid to C for rent of business premises. Compute the book profit of the firm.\pts{14}

\centerline{   \textbf{OR}}

What do you mean by Minimum Alternate Tax (MAT)? Explain when a company will have to pay MAT.\pts{14}

\medskip\rule{\linewidth}{0.2pt}




\noindent   \textbf{Question 4.}
Explain the provisions of Tax Deducted at Source (TDS) under the head ``SALARIES''. Also explain the consequences if tax is not deducted at source.\pts{14}

\centerline{   \textbf{OR}}

Mr. Rai estimated to earn the following income during the financial year 2015-16:

\begin{parts}
  \item Taxable income from House property: Rs.\ 75,000
  \item Taxable income from Business: Rs.\ 6,45,000
  \item Dividend from cooperative society: Rs.\ 16,000
\end{parts}
Determine the amount of installments payable as advance tax during the financial year 2015-16.\pts{14}

\medskip\rule{\linewidth}{0.2pt}




\noindent   \textbf{Question 5.}
Briefly describe the procedure of an appeal to the Commissioner (Appeals).\pts{14}

\centerline{   \textbf{OR}}

What are the different types of assessment? Explain any one of them.\pts{14}

\vfill
\rule{\linewidth}{0.4pt}

\begin{center}\small
\href{https://bhu-exam.vercel.app/}{Click here to visit our website}\\[4pt]
\href{https://me-aryan.vercel.app/}{Click here to visit our contributor}\\[4pt]
\href{https://quashan-me.vercel.app/}{Click here to visit our contributor}
\end{center}

\end{document}
